\section{Conclusion}\label{sec:conclusion}
Neural Bellman Operators separate the learning of continuation values from feasible economic improvement. Their useful object is a continuation that is shared across decisions, not a small residual considered in isolation. The evaluation and policy-error results specify how approximation enters decisions; the capital construction places their consequences in economic payoff units.

The signed occupation account provides a positive continuous-economy assessment of the returned NBO policies where the original norm-based sufficient bound was negative. Fresh prospective confirmation also supports the specified NBO--HJB payoff comparison. Neither result, by itself, identifies a unique incremental contribution of the learned critic relative to cached Raw. The latter question is made explicit by the common-continuation, bias--variance, and economic decision-value identities. They give a constructive positive mechanism and its exact approximation price, rather than an unconditional ranking of method names.

The numerical record retains the unresolved Raw and direct-policy comparisons, Raw's lower original work, broad high-dimensional regret bounds, and classical scalar and state-specific reversals. The completed strengthened-HJB comparison shows why the comparator specification matters: most of the enlarged-family payoff contrasts remain unresolved despite positive NBO gains. The five-method continuation experiment gives a more discriminating account of reuse. NBO has material advantages over fitted alternatives in the difficult long fifty-dimensional case and achieves practical payoff equivalence with cached sample-average optimization at lower construction and query work in six of eight final-stage cases. The sample-average method nevertheless has material payoff advantages in both long-continuation cases. All stages and all economic designs remain visible. Eighty of 128 conditional scalar candidates meet the declared accuracy margin over their entire action segments; the other forty-eight do not, and no full-vector conclusion is substituted for these conditional statements.

The original recursive-utility, preference-formation, temporal-self, and game applications remain part of the theory under their stated assumptions. The current supplement retains their complete derivations. Together, these results make the NBO program a study of learned continuation, feasible decisions, and their economic value, with policy verification and comparative method evidence serving distinct roles.
